%% kapa: exam=1 round=31 year=2020 period=1 form=A subject=economics subject_ko=경제학원론 seq=2 number=42
%% source: Q-Net 감정평가사 기출문제 (cst00309) artlSeq=5207219 file=extracted/1008/5207219/A형/A형/2020년도 제31회 감정평가사 1차 시험_1교시_A형.pdf
%% conversion: kordoc 4.24.51 -> tools/gen1.py (auto); review=proofread-sonnet; flags=none
\begin{question}{42}
A기업의 총비용곡선이 $TC=100+Q^{2}$ 일 때, 옳은 것은? (단, $Q$ 는 생산량)
\begin{choices}
  \item 평균가변비용곡선은 U자 모양을 갖는다.
  \item 평균고정비용곡선은 수직선이다.
  \item 한계비용곡선은 수평선이다.
  \item 생산량이 10 일 때 평균비용과 한계비용이 같다.
  \item 평균비용의 최솟값은 10 이다.
\end{choices}
\end{question}
\answer{4}
